\documentclass[11pt]{article}
\usepackage[T1]{fontenc}
\usepackage[utf8]{inputenc}
\usepackage{lmodern}
\usepackage[french]{babel}
\usepackage[margin=1in]{geometry}
\usepackage{enumitem}

\begin{document}

\section*{Cas d'utilisation : Créer rapport}

\begin{description}[leftmargin=4.2cm, style=nextline]
    \item[Titre de cas d'utilisation] Créer rapport

    \item[Acteurs principaux] Médecin, Administrateur (hérite les fonctionnalités du Médecin)

    \item[Description] Permet au médecin de créer un rapport radiologique par dictée vocale
          (navigateur ou mobile via QR code) ou import audio. Le système transcrit automatiquement
          via notre modèle de reconnaissance vocale médicale, normalise le texte,
          puis propose des corrections avec MedCorrect.

    \item[Préconditions]
    \begin{itemize}[leftmargin=*, nosep]
        \item Le médecin est authentifié dans le système.
        \item Le modèle de transcription est chargé en mémoire.
    \end{itemize}

    \item[Scénario principal]
    \begin{enumerate}[leftmargin=*, topsep=0pt, itemsep=0pt]
        \item Le médecin clique sur «~Nouveau rapport~».
        \item Le système affiche le formulaire de création.
        \item Le médecin saisit l'identifiant de l'examen et sélectionne la catégorie.
        \item Le système valide le format de l'identifiant et vérifie qu'il n'est pas déjà utilisé.
        \item Le médecin choisit la méthode d'enregistrement parmi :
        \begin{itemize}[leftmargin=*, nosep]
            \item Enregistrement direct via le microphone du navigateur.
            \item Enregistrement via une page web mobile après scan d'un QR code.
            \item Importation d'un fichier audio existant.
        \end{itemize}
        \item Le médecin effectue ou fournit l'enregistrement audio.
        \item Le système envoie l'audio au module de transcription.
        \item Le système applique le pipeline de normalisation et de ponctuation médicale
              sur le texte brut.
        \item Le système sauvegarde automatiquement le rapport avec le statut \textbf{Brouillon}.
        \item Le système affiche le texte transcrit et normalisé dans l'éditeur, structuré en
              sections~: \textbf{Indication}, \textbf{Résultat} et \textbf{Conclusion}
              (la section \textbf{Technique} est ajoutée pour les catégories IRM et Scanner).
        \item Le médecin vérifie et corrige si nécessaire le contenu du rapport, assisté par
              \textbf{MedCorrect} qui propose des suggestions de corrections.
        \item Le médecin choisit l'action de sauvegarde parmi~:
              \textbf{Brouillon}, \textbf{Valider} et \textbf{Enregistrer}.
        \item Le système enregistre le rapport et redirige le médecin vers son historique.
    \end{enumerate}

    \item[Scénarios alternatifs]
    \textbf{Scénario alternatif 1 : Enregistrement via page web mobile}
    \begin{enumerate}[leftmargin=*, topsep=0pt, itemsep=0pt]
        \item Le médecin sélectionne l'option «~Enregistrer via smartphone~».
        \item Le système génère une session d'appairage et affiche un QR code.
        \item Le médecin scanne le QR code avec son smartphone~; une page web s'ouvre.
        \item L'interface desktop confirme la connexion du mobile.
        \item Le médecin enregistre l'audio depuis son appareil mobile.
        \item Une fois l'enregistrement terminé, l'audio est transmis au desktop.
        \item Le desktop lance la transcription comme décrit aux étapes 7 à 9 du scénario principal.
    \end{enumerate}

    \textbf{Scénario alternatif 2 : Importation d'un fichier audio}
    \begin{enumerate}[leftmargin=*, topsep=0pt, itemsep=0pt]
        \item Le médecin sélectionne l'option «~Importer un fichier audio~».
        \item Le médecin choisit un fichier audio depuis son poste.
        \item Le système lance la transcription comme décrit aux étapes 7 à 9 du scénario principal.
    \end{enumerate}

    \item[Postconditions]
    \begin{itemize}[leftmargin=*, nosep]
        \item Le rapport est enregistré dans la base de données avec son statut.
        \item Si le statut est \textbf{Enregistrer}, une ligne d'archive est ajoutée au fichier CSV global.
        \item Le rapport est consultable depuis la liste des rapports du médecin (historique).
        \item Les administrateurs voient les rapports au statut \textbf{Valider} ou \textbf{Enregistrer}
              dans leurs propres listes.
    \end{itemize}

    \item[Exceptions]
    \textbf{Exception 1 : Format d'identifiant invalide}
    \begin{itemize}[leftmargin=*, nosep]
        \item Si l'identifiant saisi est invalide (non numérique, trop court ou année incorrecte),
              le système affiche un message d'erreur et bloque le passage à l'étape suivante.
    \end{itemize}
    \textbf{Exception 2 : Identifiant d'examen déjà utilisé}
    \begin{itemize}[leftmargin=*, nosep]
        \item Si un rapport portant le même identifiant d'examen existe déjà,
              le système refuse la création et affiche un message d'erreur.
    \end{itemize}
    \textbf{Exception 3 : Échec de la transcription}
    \begin{itemize}[leftmargin=*, nosep]
        \item Si le module de transcription échoue (audio corrompu, format non supporté ou erreur serveur),
              le système affiche un message d'erreur et propose de réessayer sans quitter la page.
    \end{itemize}
    \textbf{Exception 4 : Expiration ou perte de session mobile}
    \begin{itemize}[leftmargin=*, nosep]
        \item Si la session QR code expire ou si la connexion avec la page mobile est interrompue,
              le système notifie le médecin. Le médecin doit créer une nouvelle session
              pour relancer l'appairage mobile.
    \end{itemize}
\end{description}

\end{document}
